% GENERATED, edit the French unit or the translation table
\mxfront{expertes-front}
\mxcover{FURTHER MATHEMATICS}{Terminale, general track · option}{Course and practice booklet}{Complex numbers · Arithmetic · Graphs and matrices}{2026-2027}
\frontmatter
\chapter*{Using this booklet}
\addcontentsline{toc}{chapter}{Using this booklet}
This booklet covers the further mathematics option (mathématiques expertes) of the general Terminale. It is designed to be studied alongside the mathematics specialty course. Every chapter follows the same order.
\begin{enumerate}
\item \textbf{Fluency}: a few short questions that bring back the prerequisites, done without a calculator at the start of a lesson.
\item \textbf{Introductory activity}: an investigation carried out before the lesson. We compute, observe and state a conjecture; the proof comes afterwards.
\item \textbf{Course notes}: definitions, properties, proofs and examples. This is the written record to learn.
\item \textbf{Exercises}: in increasing order of difficulty. The first ones check a technique, the last ones require a complete written argument.
\item \textbf{Problem}: a longer study that extends the chapter.
\end{enumerate}
\begin{keypoint}[Experimenting is not proving]
A computation on a few examples, a figure or a table of values gives a conjecture. Only a proof establishes that it holds in every case. Every activity ends by saying so, as does every interactive figure in the slide decks.
\end{keypoint}
\textbf{Syllabus followed.} Syllabus of the optional further mathematics course of the general Terminale, order of 19 July 2019, special BO n\textsuperscript{o} 8 of 25 July 2019 (NOR MENE1921264A). Every proof that this syllabus lists as required is written out in the course notes. Anything beyond the syllabus is placed in a ``Beyond the syllabus'' box.

\textbf{Organisation.} The chapters are grouped by theme: complex numbers (chapters 1 to 5), arithmetic (chapters 6 to 8), graphs and matrices (chapters 9 to 11). In class, it is advisable to alternate the themes, for example a chapter on complex numbers followed by a chapter on arithmetic.
\ifmxteacher
\mxnewpage
\chapter*{Preparation guide}
\addcontentsline{toc}{chapter}{Preparation guide}
\begin{teachernote}[Content and editions]
The class option \texttt{edition=teacher} produces the teacher edition, with the solutions of the fluency questions, the activity, the exercises and the problem, together with teaching notes. The option \texttt{edition=student} produces the student edition, without solutions. The sources contain both editions.
\end{teachernote}
\paragraph{Suggested progression.} With three hours a week, one possible alternation is: C1, C6, C2, C7, C3, C8, C4, C5, then C9, C10, C11. The chapters on graphs and matrices use the conditional probability of the specialty course; place them after it has been studied there. Chapter 3 uses the scalar product of Première; the linearisation and integration exercise of chapter 3 assumes the integral calculus of the specialty course and can be postponed.
\paragraph{Durations.} The durations of the activities are estimates of actual working time, not measured with students. Plan individual reading, research in pairs, then a class discussion before the written record.
\paragraph{Required proofs.} Conjugate of a product, of an inverse, of a power; binomial formula; $\abs{z}^2=z\overline{z}$; modulus of a product and of a power; one addition formula; factorisations of $z^n-a^n$ and of $P(z)$ by $z-a$; number of roots at most equal to the degree; determination of $\mathbb{U}_n$; writing the gcd in the form $au+bv$; Gauss's theorem; infinitude of the primes; number of walks of length $n$ and probability of transition in $n$ steps. They appear in the course notes of the corresponding chapters.
\paragraph{Problems.} The end-of-chapter problems take up ``possible problems'' cited by the syllabus. They are not required: they serve for independent work, homework or extension.
\fi
\tableofcontents
\mainmatter
